\section{Start with the requested object}
A statement's story is less useful than its output and constraints. Write
one sentence saying what counts as a feasible answer. Then identify the
quantity optimized or counted. ``Find a route'' can mean any route, fewest
edges, smallest weight, a simple path, or a route that obeys a time limit.
Those lead to different algorithms even on the same graph.

Use the code references after this choice. The math book holds formulas and
numeric contracts; the implementation book holds foundational helpers and
KACTL specialists. This book supplies reductions, traces, code adaptations,
and failure cases. The cases marked synthetic isolate one decision; the
named contest cases link their sources. None of the example input/output
wrappers has been submitted to a judge. The five new helper modules have
local independent-oracle tests.

\subsection{A five-line design before typing}
\begin{enumerate}
\item \textbf{Object.} A set, subsequence, contiguous interval, path, walk,
assignment, or number? Preserve this distinction in the representation.
\item \textbf{Information.} What must be known to decide the next legal
move? If two histories can have different legal continuations, they cannot
share one DP or search state.
\item \textbf{Naive method.} Enumerate the choices or write the recurrence.
Identify the repeated work rather than guessing an advanced structure.
\item \textbf{Reduction.} Explain why discarded information cannot change
the answer. For greedy, explain why an optimal solution can use the choice.
For binary search, prove feasibility stays true as the bound increases.
\item \textbf{Cost.} Count states, transitions, stored values, and numeric
bounds. Multiply dimensions; a small-looking state with three indices can
be too large. Compare the estimate with the actual time/memory limits.
\end{enumerate}

\subsection{Terms used in the graph cases}
A vertex represents a state; an edge represents one permitted transition.
An adjacency list stores each vertex's outgoing neighbors. Directed edges
allow travel only in their stated direction; undirected edges are inserted
both ways. A walk may repeat vertices, while a simple path may not. BFS
uses a first-in, first-out queue, so it explores equal-cost transitions by
increasing number of moves. A DP state stores an answer for a subproblem;
its transition combines already solved states in a dependency-respecting order.

\section{Recognition routes}
A cue suggests a candidate, not a proof. Read the condition in the right
column before selecting the tool. The page references lead to derivations.

\subsection{Numbers and counting}
\noindent\begin{tabular}{@{}p{.28\columnwidth}p{.60\columnwidth}@{}}
\textbf{Clue} & \textbf{Model and required check}\\\hline
Huge integer, allowed error & Keep significant digits and prove relative error. Exact output would invalidate this. Fuzzy, p.~\pageref{case:fuzzy}.\\
Exact decimal prices & Parse scaled integers with a common scale; gcd or divisibility then applies exactly. Scaling, p.~\pageref{case:orbits}.\\
Unknown digits and divisibility & State is position plus prefix remainder; add any other future-relevant property. Enigma, p.~\pageref{case:enigma}.\\
Repeated process on a circle & Work modulo the size. Repeated step $d$ visits $n/\gcd(n,d)$ positions; there are $\gcd(n,d)$ orbits.\\
Count pairs or expected total & Fix one contribution, count how often it appears, then sum. Specify ordered/unordered and replacement first.\\
Synchronize periodic schedules & Solve congruences. CRT with shared factors requires residue differences divisible by the gcd.\\
Smallest missing subset sum & Sorted positive values may cover a whole interval. Prove coverage before skipping enumeration, p.~\pageref{case:missing}.\\
\end{tabular}

For a contribution example, array element $a_i$ appears in $(i+1)(n-i)$
subarrays: choose a left endpoint from $0..i$ and right endpoint from
$i..n-1$. Therefore the sum of all subarray sums is
$\sum_i a_i(i+1)(n-i)$. This avoids enumerating the subarrays. It does not
count subarrays whose sum equals a target; that predicate couples elements
and needs a different reduction, p.~\pageref{case:exact-sum}.

\subsection{Graphs and dependencies}
\noindent\begin{tabular}{@{}p{.28\columnwidth}p{.60\columnwidth}@{}}
\textbf{Clue} & \textbf{Model and required check}\\\hline
Fewest equal-cost moves & BFS on complete states, possibly position plus mask/parity. Product states, p.~\pageref{case:state-bfs}.\\
Costs only zero and one & 0--1 BFS; discovery is not final until distance improvements stop. Arbitrary costs need another algorithm.\\
Minimum nonnegative total cost & Dijkstra; add resource/phase information to vertices when future choices depend on it.\\
Can events be ordered? & Directed edges from prerequisite to dependent; topological order detects cycles and supports DAG DP.\\
Connect everything cheaply & MST, usually Kruskal with DSU; this minimizes edge total, not source-to-target distance.\\
Mutual directed reachability & SCC, then a condensation DAG if processing components. Undirected components are not a substitute.\\
Two groups of objects & Bipartite matching for one-to-one assignment; flow for quotas. Define each edge's meaning, p.~\pageref{case:matching}.\\
Avoid spreading hazards & Earliest hazard arrivals become deadlines for a constrained player BFS, p.~\pageref{case:hazards}.\\
\end{tabular}

A tree is a connected undirected graph without cycles. Rooting gives each
nonroot vertex one parent. Process parents before children for depth/root
sums; reverse that order for child-to-parent DP. LCA answers path ancestry,
not arbitrary range aggregates. For path sums with only static edge weights,
root distances give $d(u,v)=D[u]+D[v]-2D[\mathrm{LCA}(u,v)]$; dynamic path
updates need more structure. The supplied root order is not a contiguous
subtree Euler tour.

\subsection{Arrays, ranges, and optimization}
\noindent\begin{tabular}{@{}p{.28\columnwidth}p{.60\columnwidth}@{}}
\textbf{Clue} & \textbf{Model and required check}\\\hline
Static interval sum & Prefix sums and subtraction; negative values are fine.\\
Many range additions, one final array & Difference endpoints, then prefix scan. Interleaved queries need another structure.\\
Point changes, range sums & Fenwick; use delta for assignment. Prefix order statistics require nonnegative frequencies.\\
Associative range summary & Segment tree with a correct identity and merge; sum, min, max and ordered concatenation differ.\\
Overlapping coverage & Sort intervals, merge, sum lengths. Depth of coverage instead needs events, p.~\pageref{case:intervals}.\\
Minimize largest feasible bound & Prove a monotone decision predicate, then binary search; its internal solver still needs proof, p.~\pageref{case:capacity}.\\
Best earlier compatible choice & Write $dp[i]=w_i+\max dp[j]$. A dominance/prefix query may replace the scan, p.~\pageref{case:fundraising}.\\
Small integer budget, optional items & DP by processed items and capacity; scan direction controls reuse, p.~\pageref{case:knapsack}.\\
Few independent binary actions & Enumerate subsets if the operation count fits. Commuting toggles give a concrete example, p.~\pageref{case:flip}.\\
\end{tabular}

\section{A nearby problem that needs a different tool}
\label{case:exact-sum}
\subsection{Signed exact sums: replace the window by prefix frequencies}
Synthetic task: count nonempty contiguous subarrays with sum $K$. If all
values are nonnegative, expanding a window cannot decrease its sum and
removing the leftmost value cannot increase it. This supports the window
helper in the implementation book. Negative values break those implications:
for $[3,-2]$ and $K=1$, discarding the first 3 when it looks too large loses
the valid two-element subarray.

Define $P[0]=0$ and $P[r]=a_0+\cdots+a_{r-1}$. An interval $[l,r)$ has
sum $K$ exactly when $P[l]=P[r]-K$. Scan $r$ left to right and count earlier
prefix values. Initialize frequency of 0 to one. Query before inserting
$P[r]$, so $l<r$ and the empty subarray is excluded.

Trace $[1,-1,1]$, $K=1$: start with frequencies $\{0:1\}$. Prefix 1 finds
one earlier 0; prefix 0 finds no earlier $-1$; prefix 1 now finds two earlier
zeros. Total 3, from $[0,1)$, $[0,3)$ and $[2,3)$. Storing only whether a
prefix exists would lose multiplicity.

The module uses an ordered map for guaranteed $O(n\log n)$ time and $O(n)$
space. A hash map changes this to expected linear time, with hash/runtime
tradeoffs. Prefix sums and $P[r]-K$ must fit the numeric type. For minimum
length instead of count, keep the latest matching index; for maximum length,
keep the earliest. For arbitrary intervals with inequalities, an equality
map alone is insufficient; consider ordered prefix queries.

\subsection{Shared code for this book}
The new modules use namespace \texttt{finals\_cases}. Copy their shared
prelude and the listed dependencies once. \texttt{Fundraising.h} additionally
uses the implementation book's segment tree; \texttt{CapacitySearch.h} uses
its binary search. The includes below point to the repository files; when
typing into one contest file, paste those dependencies instead of retaining
local paths. Existing helpers keep \texttt{finals\_core} or
\texttt{finals\_math}; qualify a call or use the corresponding namespace.
\sourceimport{finals-patterns}{Prelude.h}
\sourceimport{finals-patterns}{ExactSum.h}

Practice variant: count zero-sum nonempty intervals in $[0,0]$. Answer: 3.
The three prefixes are all zero; each new prefix pairs with every earlier
one. Inserting before querying would incorrectly count an empty interval
at each step.
